不完全性定理的证明要求我们能够用理论本身的语言——即算术语言——来谈论理论（如 $\Th{PA}$）内部的可证性。但算术语言只涉及数，不涉及公式或推导。解决这个问题的办法是定义一个从公式和推导到数的系统映射。与公式或推导相关联的数称为它的\textbf{G\"odel 数}。如果 $!A$ 是公式，$\Gn{!A}$ 就是它的 G\"odel 数。我们证明了，公式上的重要操作会转化为关于相应 G\"odel 数的原始递归函数。例如，$\Subst{!A}{t}{x}$，即用项~$t$ 替换~$!A$ 中变元~$x$ 的每一自由出现这一操作，对应于算术函数 $\fn{subst}(n,m,k)$，将其应用于 $!A$、$t$ 和 $x$ 的 G\"odel 数时，便得出 $\Subst{!A}{t}{x}$ 的 G\"odel 数。换言之，$\fn{subst}(\Gn{!A}, \Gn{t}, \Gn{x}) = \Gn{\Subst{!A}{t}{x}}$。同样地，推导的性质也转化为关于相应 G\"odel 数的原始递归关系。具体而言，性质 $\fn{Deriv}(n)$（当 $n$ 是自然演绎中某个正确推导的 G\"odel 数时成立）是原始递归的。证明它们是原始递归的，需要相当多的工作，有时还需要一些巧思，并且本质上依赖于“对序列的操作是原始递归的”这一事实。如果理论~$\Th{T}$ 是可判定的，那么我们就能用 $\fn{Deriv}$ 来定义一个可判定的关系 $\Prf[\Th{T}](n, m)$，当 $n$ 是从~$\Th{T}$ 出发对 G\"odel 数为~$m$ 的语句所作推导的 G\"odel 数时，该关系成立。若~$\Th{T}$ 的公理集是原始递归的，则该关系也是原始递归的；若~$\Th{T}$ 的公理只是可判定的但非原始递归的，则该关系只是一般递归的。